% ============================================================
%  R. Nielsen, Svar til Hr. Dr. phil. Pastor Zeuthen i Anledning af hans
%  Sendebrev om Tro og Viden
%  [Reply to Pastor Zeuthen, Dr. phil., on the Occasion of His Open
%  Letter on Faith and Knowledge]
%  Copenhagen: Gyldendal (F. Hegel), printed by J. H. Schultz, 1866. 16 pp.
%  TRANSLATION (English)
%  Source of truth: transcription.tex (Danish), COMPLETE.
%  Mirrors transcription.tex 1:1; every printed page carries \opage{N}
%  at the corresponding point in the English. See TRANSLATION-PLAYBOOK.md.
%  Translated 22 Sep 2026 in two batches (pp. 1--8, 9--16).
%  Latin and Greek kept as printed (the Greek of p. 6 is unaccented in the
%  print and is copied so). Danish work titles cited by Nielsen are left in
%  Danish; Martensen's "Den christelige Dogmatik", which Nielsen quotes
%  throughout in quotation marks, is rendered ``the Christian Dogmatics''.
%  Printer's errors in the Danish (Natuvidenskaben, traditonelle, ...) are
%  not reproduced in the English; each is noted in the transcription.
%  Bible passages are translated from Nielsen's Danish, in the register of
%  the Authorized Version where the wording allows.
%  Terms: Tro = faith; Viden = knowledge; Videnskab = science;
%  Erkjendelse = cognition; Under = wonder; Mirakel = miracle;
%  det Underfulde = the wonderful; Naturlære = doctrine of nature;
%  Ideelære = doctrine of ideas; Ideevidenskab = science of ideas;
%  Mirakelvidenskab = science of miracles; Sendebrev = open letter;
%  Villie = will; Opholdelse = preservation; Afgrund = abyss.
% ============================================================
\documentclass[12pt,a4paper]{article}
\usepackage[utf8]{inputenc}
\usepackage[T1]{fontenc}
\usepackage{libertinus}
\usepackage{libertinust1math}
\usepackage{textalpha}   % Greek typed directly (p. 6)
\usepackage[english]{babel}
\usepackage[protrusion=true,expansion=true]{microtype}
\usepackage{setspace}
\usepackage[margin=1.3in]{geometry}
\usepackage{fancyhdr}
\usepackage{marginnote}
\newcommand{\opage}[1]{\marginnote{\footnotesize\textit{[#1]}}}
\usepackage{hyperref}

\hypersetup{
  pdftitle={Reply to Pastor Zeuthen, Dr. phil., on the Occasion of His Open Letter on Faith and Knowledge},
  pdfauthor={R. Nielsen},
  hidelinks
}

\setlength{\headheight}{15pt}
\pagestyle{fancy}
\fancyhf{}
\rhead{Nielsen, \emph{Reply to Pastor Zeuthen} (1866)}
\cfoot{\thepage}
\renewcommand{\thefootnote}{\ifcase\value{footnote}\or *)\or **)\or ***)\else
  \arabic{footnote})\fi}
\setlength{\parindent}{1.5em}
\setstretch{1.3}

\begin{document}

% --- p. 1 ---
\opage{1}%
% Title page, mirroring the Danish (transcription.tex p. [1]).
\begin{center}
{\Huge Reply}\\[1.5em]
{\small to}\\[1.5em]
{\LARGE Pastor Zeuthen, \textit{Dr. phil.},}\\[1em]
{\large on the Occasion of His Open Letter}\\[1em]
{\Huge\bfseries on Faith and Knowledge.}\\[1em]
{\small By}\\[0.5em]
{\large\bfseries R. Nielsen.}\\[6em]
{\bfseries Copenhagen.}\\[0.5em]
Published by the Gyldendal Bookshop (F. Hegel).\\[0.3em]
{\footnotesize Printed by J. H. \emph{Schultz}.}\\[0.5em]
1866.
\end{center}
\clearpage
% --- p. 2 ---
\opage{2}%
\begin{center}
{\large\bfseries Most esteemed Doctor!}
\end{center}

We could, so far as I can see, come to agreement on very
essential points, if only we could agree that a science of
miracles, a science that explains miracles, belongs among the
impossible things. You ask me in your Open Letter (p.~7) how one
can declare it ``a settled matter that science denies the wonder,''
and you add straight afterwards: ``As far as natural science is
concerned, one must indeed grant that its concept of nature is not
theological, and that it has no place for what we in the religious
sense call the wonderful; but is natural science all science, the
highest science?'' I answer: that besides natural science there is
indeed a science of ideas; but of this latter, in my view, three
things hold: 1)~it must be built on the mathematical-empirical
foundation of the doctrine of nature; 2)~it must, just as much as
the doctrine of nature, be a cognition by virtue of laws; it can
3)~not be thought to stand in any essential or principled conflict
with the doctrine of nature. If, now, the doctrine of nature as
science---as science, mark you---must deny the possibility that
``one really dead'' can really rise from the grave (so that no
natural-scientific account whatsoever can be given of the miracle
of the resurrection), then the doctrine of ideas, as surely as it
would count as science, must deny the same. For the doctrine of
nature would be false if, in its interpretation of the facts, it
declared impossible what was nonetheless possible according to a
scientific conception of the idea of nature; conversely,
% --- p. 3 ---
\opage{3}the doctrine of ideas would be only hollow speculation if, by its
extension of the possibilities, it came into conflict with the
first inviolable principles of the doctrine of nature. To try to
bring about a reconciliation between religion and science by
sowing discord between science and science will in no way do.

The old dream---for I can now regard it as nothing other than a
dream---of a highest science that should be able to find out the
``thinkability'' and ``possibility'' of the miracle rests, if I am
not mistaken, on a couple of ambiguities: the ambiguity in the word
\emph{wonder}, and the ambiguity in the use theology makes of the
teleological.

What is the wonder, when theology sets about treating it
``objectively,'' other than a great ambiguity? ``Everything''---says
an acute theologian---``is a wonder. Every event is the effect of
the Godhead, grounded immediately in the divine power; but it shows
itself as an effect of nature grounded mediately in God. In the
objective sense, then, everything is wonder or nothing is wonder;
the awakening of the seed to the life of spring is a wonder as much
as the awakening of the dead from the grave, healing by medicines as
much as by the laying on of hands, and every point of support for
the concept falls away. Nature's system of forces is a deep that is
eternally unfathomable, and we only shame ourselves if we presume to
draw boundary lines where there are none.'' But that a doctrine of
ambiguity which places ``the awakening of the dead from the grave''
and ``the awakening of the seed to the life of spring'' on the same
step in the series of possibilities is not exactly fitted to
reconcile faith and knowledge, but ought rather to be assigned ``a
place of its own outside science,'' most people in our time will
surely concede.

To the ambiguous use of the word wonder there corresponds in
theology an equally ambiguous use of the words \emph{teleology} and
the \emph{teleological}. For there is a difference, indeed a very
essential difference, between a teleology that is the higher
expression of the fulfilment of laws, a presentation of nature's
real ends, and a fantastic teleology that soars above the laws of
science. Since, then, it belongs in my opinion to the first
requirements of philosophical culture to
% --- p. 4 ---
\opage{4}be oriented regarding the difference between what is determined by
law and what is fantastic, I long ago treated the ambiguities
belonging here in the Propaedeutics, and I set down here, by way of
example, a small specimen: ``A. Do you suppose that the natural
scientist's concept of the law of nature stands in an essential or
in a merely accidental relation to natural science? C. In an
essential one, naturally! A. You quoted above a passage from
H.~C. Ørsted's \emph{The Spirit in Nature}, and will then also
remember that in the same work there is an essay, `The Relation of
Natural Science to Several Important Objects of Religion,' whose
first part, headed `The Immutability of the Laws of Nature,' begins
thus: `In my book, The Spirit in Nature, I have with full conviction
embraced the old doctrine of the immutability of the laws of
nature\ldots.' But if the laws of nature are immutable, how are we
then to interpret sentences such as these: `The awakening of the
seed to the life of spring is a wonder, as much as the awakening of
the dead from the grave, healing by medicines as much as by the
laying on of hands'? C. The hermeneutic key lies, surely, in the
words: `Every event is the working of the Godhead, grounded
immediately in the divine power, but it shows itself as an effect
of nature grounded mediately in God. In the objective sense, then,
everything is wonder or nothing is wonder.' A. So if the dead
awoke in the grave every spring, when the seed `awakens to life,'
the former would be in just as good agreement with `the
immutability of the laws of nature' as the latter? C. Yes, why not,
if such an awakening really were \emph{teleological}? For nature
is''---here a sentence is borrowed from \textit{Dr.} H.~Martensen's
``Christian Dogmatics''---``a system that finds itself in a
continued teleological development, a continued creation!'' Why
should the dead not be allowed to awaken every spring, provided
only that they awoke teleologically, in accordance with a continued
creation? A. That arsenic sulphide can turn into water, water into
wine, wine into coal, and coal into gold, is then in perfect
agreement with the immutability of the laws of nature? C. Provided,
mark you, that it happens \emph{teleologically}, in accordance
with a continued creation! A. That a
% --- p. 5 ---
\opage{5}man walks on the sea without sinking is then just as consonant
with the immutability of the laws of nature as that someone falls
into the water and drowns? C. Man, who is otherwise subject to the
laws of gravity, can very well walk on the sea when it happens
\emph{teleologically}, in accordance with a continued creation! A.
But then, as far as the immutability of the laws of nature is
concerned, there is presumably nothing offensive to deeper thought,
when we rightly consider it, in a Pope, Leo XII, having been able to
beatify the Minorite Julianus, who in our unbelieving, revolutionary
nineteenth century made roasted birds fly? C. That I for my part
take the wonder of the roasted birds for a stupid fable, a priestly
lie, does not arise at all from my seeing in the event itself any
breach of the immutability of the laws of nature, but simply from my
not being a Catholic and consequently being unable to convince
myself that the event has taken place \emph{teleologically}, in
accordance with a \emph{continued creation}. Were I a Catholic and
had a Catholic's faith, I should surely be able to reconcile my
faith with my knowledge, my concept of creation with my concept of
nature; I should then surely, in consideration---indeed in dogmatic
consideration---of the teleology in the roasted birds, be able to
explain the miracle thus: ``Nature''---here again a sentence is
borrowed from \textit{Dr.} H.~Martensen's ``Christian
Dogmatics''---``does not deny the concept of creation, and the
wonder is precisely the point at which nature's dependence on the
creating freedom becomes perfectly manifest.'' If the eternal Reason,
the triune God, the absolute Idea, wished to give us a sign that
could show the unbelief of these times in a really conspicuous way
that nature, a system that finds itself in a continued
\emph{teleological} development, does not deny the concept of
creation: how could the eternal Reason, working according to
\emph{immutable laws of nature}, then choose a sign more
\emph{teleological} than precisely this sign of the Minorite's
roasted birds? For this stands dogmatically firm: if the birds are
roasted \emph{teleologically}, in accordance with a process of
nature continued in creation, then they must also be able to fly
\emph{teleologically}, in accordance with a process of creation
continued in nature. A. You are a true theologian; you can
reconcile faith
% --- p. 6 ---
\opage{6}%
and knowledge.'' (Lectures on ``Phil.\ Propæd.'' 1861--62;
pp.~158--59).

In this short exchange, Doctor, you will surely find neither
distortion nor exaggeration. It is, after all, a well-known matter
that the Catholic scholastic, leaning on the authority of the
Church, set firmly on his medieval traditional foundation, has an
altogether different assurance in explaining and defending his
saints' miracles than the Protestant theologian, dazzled by the
enlightenment of reason and entangled in seventeen half-scientific
considerations, has in explaining the works either of Christ or of
his Apostles.

To convince ourselves still further of what it comes to in reality
when the supposedly higher science sets about explaining and
grounding the wonders of revelation, we need only cast a glance at
``the Christian Dogmatics''; for we have, happily, in our recent
literature the dogmatics that can with reason be said to represent a
summit of theological scientificity.

Yes, if any science should be able to explain the wonder, it would
certainly have to be dogmatics. What, then, has our ``Christian
Dogmatics'' brought forth toward a scientific understanding of the
three fundamental wonders of revelation: \emph{Creation, the
Incarnation} and \emph{Inspiration}?

Of \emph{Creation} it is taught that God ``brings forth a creature
that at the same time brings forth and develops itself in given
independence.'' If it is to be possible to connect any meaning with
so obscure a sentence, the meaning must evidently be this: that the
activity of the Creator is at the same time an activity of nature,
that the wonder at the same time takes place in a natural way. With
this it then certainly agrees that the world is created on the one
side with necessity, on the other with freedom; on the one side out
of nothing, on the other out of something; that creation is at
bottom cosmogony, and cosmogony creation. ``Judaism,'' it is said,
``conceives the world predominantly as \textit{creatura}, not as
\textit{natura}, as κτισις, not as φυσις. But just for that reason
the whole significance of creation is not disclosed to Judaism'';
and so on. Once the miracle has in this way become
% --- p. 7 ---
\opage{7}%
an event of nature and the event of nature a miracle, once nature
and creation have been lumped together so that nature is to a
certain degree created and the created is nature, then it becomes an
easy matter to introduce gradations, not only in the ``given
independence'' of the created, but in the miracle of creation
itself; the one miracle of creation then becomes greater, stronger
and more remarkable than the other: we get a dogmatic thermometer
for miracles. Thus ``the Christian Dogmatics'' sees in Christ at
once ``the completed creation of human nature and the God become man
himself,'' sees, in other words, how the Uncreated is himself
created: this, now, is the very greatest miracle. The creation of
``the first Adam'' is indeed also a great miracle. ``The first Adam
is \emph{created} in a sense in which none of his descendants is'';
but this great miracle of creation nevertheless stands many degrees
below the miracle of the Incarnation. Now, just as the miracle of
the creation of the first Adam stands many speculative degrees below
the miracle of the second, so the miracle of the creation of Eve,
though in itself a respectable miracle, stands at least two and a
half degrees below the miracle of the creation of Adam. So many
degrees in the miracle of creation, so many meanings of the word
creation! If the second Adam is created in a sense in which no other
creature is; if the first Adam is created in a degree and sense in
which none of his descendants is; if Eve is created in a degree and
sense in which no other woman is; the heroes and the geniuses again
created in other degrees and other senses than ordinary human
beings: then one sees that the miracle-degrees of creation must be
able to vary between great nodal points almost as continuously as
the natural degrees of heat. This fluid---if at times viscous---unity
of nature and creation then at last, by the addition of
``traducianism'' and ``creatianism,'' becomes remarkably thin. ``It
is the truth of traducianism that every human individual is an
offspring of the natural activity of the race.'' But it is ``the
truth of creatianism that this mysterious natural activity is the
organ and means of God's
% --- p. 8 ---
\opage{8}%
individualizing creative activity.'' Creation has ceased, for it has
been replaced by preservation; but preservation is a continued
creation, so creation has not ceased; creation has to a certain
degree ceased, to a certain degree not ceased. This, in sum, is the
explanation ``the Christian Dogmatics'' gives us of the wonder of
creation.

And what does ``the Christian Dogmatics'' teach about the wonder of
the Incarnation? Among other things the following: ``When it has
been found self-contradictory that the eternal Logos becomes man,
`because the Eternal and Omnipresent cannot let himself be born in
the midst of time,' then it certainly cannot be the task to think
that with the Incarnation the eternal Logos should cease to exist in
his general revelation in the world, or that the Logos as a
self-conscious personal being should be enclosed in a mother's womb,
be born as a child, grow in knowledge, and so on; for this notion
annuls the very concept of birth. Birth in time necessarily brings
with it the notion of a progression from the unconscious to the
conscious, from possibility to actuality, from seed-corn and germ to
full-grown organization. That the divine Logos lets himself be born
therefore means that he sinks himself into the bosom of humanity as
possibility, as holy seed, in order to be able to rise up in the
midst of the human race in mediating and saving
\emph{human} revelation.'' Here, then, we have the explanation. But,
Doctor, could we not easily agree that this explanation itself
stands in great need of an explanation? For what is it that the
divine Logos has properly ``sunk'' into ``the bosom of humanity'':
\emph{himself} or \emph{not himself}? If it is not himself, not his
proper self, then neither is it the Logos himself that is
incarnate. The being whose proper self cannot be ``sunk'' as
possibility into a mother's womb cannot be incarnated. But if it
really is himself, his possible, essential self, yet not the
eternally self-conscious personal self of which we are speaking
here: then the divine Logos must have a double self, one that can be
incarnated and one that cannot be incarnated. The Logos that can be
incarnated is therefore in himself, in his proper self, another than
the eternal personal Logos
% --- p. 9 ---
\opage{9}%
himself; in other words: according to the explanation of ``the
Christian Dogmatics,'' the eternal personal Logos himself has not
become incarnate at all.

And what interpretation, finally, does ``the Christian Dogmatics''
give us of the wonder of Inspiration? It is the event at the feast
of Pentecost (Acts, ch.~2, vv.~1--12). The sound from heaven, the
rushing mighty wind, the tongues of fire, the gifts of tongues, that
``they were all filled with the Holy Ghost, and began to speak with
other tongues, as the Spirit gave them utterance''---this was what
was to be explained. ``The new creation,'' says the Dogmatics,
``in no respect falls short of the old as far as the miraculous is
concerned; here `a supernatural beginning' is made.''%
% The Danish prints two opening marks („staaer ... „en overnaturlig
% Begyndelse“) and one closing mark here; the English nests the inner
% quotation and closes both.
{} But if we look more closely, then by dogmatic-scientific
measurement of degrees---for the miracle of the Spirit naturally has
its degrees and gradations too---the miracle has been placed so
peculiarly that it almost looks as though ``the process of nature''
in ``the new creation'' had gained a disquieting preponderance over
``the process of creation.'' For the state in which the Apostles
found themselves when at the feast of Pentecost they began to speak
with tongues is described as ``a state in which consciousness is
not master of itself, a state of ecstasy, of rapture, which was the
expression of the breaking through and stirring of the new spiritual
powers in the depths of the soul, and which bore the stamp of a
spiritual state of nature rather than of the clear life of
consciousness''; truly, among the many ``processes of creation'' of
sacred history, a most remarkable ``process of nature''! That this
``spiritual state of nature'' of ``ecstasy'' and ``rapture'' must
have lasted rather long and worked on a considerable scale is
especially evident from the remarkable influence it has exercised on
the author of the Acts of the Apostles. For the Evangelist Luke,
undoubtedly on account of a certain ``spiritual state of nature,''
has come in his account to stand the whole miracle on its head. What
Luke relates is, as is well known, this: that the ``devout men, out
of every nation under heaven'' heard the Apostles speak in such a
way that they ``were amazed and marvelled, saying one to another,
Behold, are not all these which speak Galilaeans? And how
% --- p. 10 ---
\opage{10}%
hear we them speak, every man in our own tongue, wherein we were
born? Parthians, and Medes, and Elamites, and we that dwell in
Mesopotamia, Judaea, and Cappadocia, Pontus and Asia, in Phrygia and
Pamphylia, the parts of Egypt and of Libya, about Cyrene, and we
Romans that dwell here, Jews and proselytes, Cretes and Arabians?
\emph{we do hear them declare the great works of God in our
tongues}.'' More clearly, more fully and more solemnly the author
could not have expressed himself, if he wished precisely to lay
weight on the wonder's showing itself just in the gift of tongues,
and the gift of tongues in this: that the simple Galileans, people
who must otherwise be supposed able to speak only one single
language, suddenly, by the prompting of the Spirit, came to speak
the many different foreign languages. The Dogmatics, which has
speculated the miracle out, knows better, however. ``The Christian
Dogmatics'' declares expressly ``that it is not necessary to think
here of different languages, but of the circumstance that `they'''
(devout men out of every nation under heaven) ``each for himself
felt addressed in his own innermost \emph{peculiarity}, that the
inner man in each was so addressed by this tongue of rapture that
the soul felt itself released from its accustomed, its natural
limitation and in a wonderful manner became manifest to itself.''
By this explanation, then, the miracle of Pentecost has literally
been turned upside down. The miracle which in ``the Acts of the
Apostles'' happens to the Apostles is explained in ``the Dogmatics''
by ``a state in which consciousness is not master of itself,''%
% The Danish opens this quotation („en Tilstand, ...) and prints no
% closing mark; it is closed here, where the phrase quoted on p. 9 ends.
{} and is to that extent to be regarded as an improper miracle;
whereas the natural joy and astonishment of the foreign ``devout
men'' at that improper miracle is transformed in ``the Christian
Dogmatics'' into the proper miracle, namely the miracle that every
one of those present, although he ``in a wonderful manner became
manifest to himself''%
% The Danish prints only the closing mark of this quotation (after
% "aabenbar"); the opening is supplied from the passage quoted above.
{}---so that here, in this respect, there could be no thought of
illusion---was nevertheless (naturally by the miracle; for such a
thing is possible only by a miracle) placed under the illusion that,
in hearing a tongue ``wherein'' he was notoriously not born, he
precisely imagined he heard ``the tongue wherein he was born.''

% --- p. 11 ---
\opage{11}%
If now these specimens of the profound ``searchings'' and objective
explanations of our ``Christian Dogmatics'' make an impression on
you, Doctor, similar to the one they make on me, then we shall soon
agree to give up the hope of the supposedly higher science, the
science of miracles explained in the Spirit.

It is the miracle and again the miracle that, in the strictest
sense, ties knots for science. We cannot untie these knots; let us
honestly admit it: there is no other way out! Yet what am I saying?
There is a way out after all, namely the way out of covering the
knots with a couple of ``valuable wisps of straw'' and pouring
nonsense over the ``wisps of straw.'' This is undeniably a way out,
as far as the treatment is concerned. But against it this must be
said: the treatment is not scientific; and the scientific was
supposed to lie in the treatment.

Science is not religious, religion not scientific; astronomy is not
devout, devotion not scientific; mathematics is not prayerful,
prayer not scientific; botany is without sacrament, the sacrament is
not scientific. But however convinced I am that faith and knowledge,
religion and science must---for the sake of a higher
reconciliation---sue for divorce, and that theology, therefore,
which as mediator would wear the cloak on both the religious and the
scientific shoulder, must go; just as convinced am I that a
revelation which has ``the Word of Life'' in itself must also have
the light in itself; that there must, in other words, be a religious
thinking, thorough, clear, convincing, a thinking which, in
accordance with its peculiar presuppositions, neither is nor can be
nor will be scientific, but which nonetheless in sharpness,
certainty, endurance and consistency is a match for science.

You remind me (Open Letter, p.~9) of Ørsted's words: ``With truth
the good is indissolubly bound, and the \emph{full} truth brings its
own comfort with it,'' adding with reason: ``it is by virtue of a
faith, a trust, and so by virtue of his subjectivity, that he says
this. But how should one come to the cognition of the full
% --- p. 12 ---
\opage{12}%
truth if one would not \emph{respect} the factual? So long as
something given in experience remains standing and testifies against
the theory one has formed, the theory cannot be true. But to what is
thus given there may also belong the human heart's \emph{need} for a
living God and for his grace, so that only in faith in it can it
find peace.'' In this I am now wholly and entirely agreed with you,
Doctor, just as we also agree that ``God as the \emph{product of
thought} cannot be the object of faith.''

But how are we to come to agreement about the remarkable relation in
which you (Open Letter, p.~11) have placed the God of faith to human
reason? It looks to me somewhat questionable. For if by reason and
reason, thinking and thinking, I am to think the same thing in both
sentences, then it is truly a difficult piece. Let me repeat what
stands here: the God of faith can, you say, \emph{to that extent}
(namely, as the God who is the product of thought cannot be the
object of faith) be called ``the true abyss for human reason''; and
immediately thereupon you add: ``But the God of faith must then''
(namely, when he is ``the true abyss for human reason'') ``be more
than this, more than him at whom all reason, all thinking must
stop''\ldots\ How, I ask, can the God of faith both be ``the
\emph{true} abyss for human reason'' and yet still more than the
true abyss, and so an abyss to such a degree that he is not an
abyss; or---let us now consider properly how ``it must be
thought''---is the God of faith perhaps an abyss that stops human
reason in such a way that the abyss itself stops before reason, just
as reason was to stop before the abyss; or is it perhaps so, that
the abyss becomes reason while reason becomes abyss; or perhaps so,
that reason first stops before the true abyss when it does not stop
at all; or perhaps\ldots\ Doctor, that piece about the abyss has a
strange aftertaste! It tastes speculatively of ``Christian
Dogmatics''; I cannot get it down.

% --- p. 13 ---
\opage{13}%
You ask (p.~12): ``Do not truth and goodness really stand in so
essential a relation to God that one cannot be said to have any
practical cognition of God at all, if one does not assume this union
of truth and goodness with God?''%
% The Danish prints no opening mark before "Staaer"; supplied here.
{} I answer: that the believer must certainly assume that truth and
goodness are united in God; but that a science of this union is, in
my view, impossible is seen directly from the opposition between the
idea of truth and the idea of goodness on which my conception of the
difference in principle between theoretical and practical, scientific
and religious cognition is founded. I by no means overlook the
interaction between truth and goodness, between the ideals of
thought and of will. ``It goes without saying,'' I have indeed said,
``that the ideals of thought and will must be able to enter into
just as many-sided an interrelation as thought and will themselves;
every attempt to draw a definitively fixed boundary line that would
separate theoretical cognition from practical is to that extent
unreasonable and meaningless. But there is here nonetheless a
boundary, a difference in principle, which arises from the fact
that in theory all the powers of the soul are subordinated to
thinking, whereas in practice they are subordinated to the will.''
(Om theoretisk og praktisk Erkjendelse, p.~10).%
% The quoted passage is rendered as in this repo's translation of that
% article (Nordisk Universitets-Tidsskrift 10, Tillægshefte, p. 72;
% Nielsen cites a separately paginated printing, p. 10).
{} You have here misunderstood my popular exposition, as though I set
the will up beside thought in an external way and, with one-sided
attention to the practical, had an eye only for the proposition
\textit{bonum est, quia deus vult}, but forgot ``to add the other as
no less justified: \textit{deus vult, quia bonum est}.'' On this
occasion, Doctor, I should not like to burden you with metaphysical
sentences, for I know from experience that such things lead
nowhere; but if at some convenient moment you should come to leaf
through Grundideernes Logik II (pp.~363--84), you would easily
convince yourself that over the proposition \textit{bonum est, quia
deus vult} I have by no means forgotten its opposite: \textit{deus
vult, quia bonum est}.

You think (Open Letter, p.~11) that I wrongly make
% --- p. 14 ---
\opage{14}%
religion practical from first to last, as though it were not ``bound
by any essential bond to a knowledge and to theoretical cognition,
as though in religion thoughts meant nothing at all in and for
themselves, but were exclusively practical and subordinate to the
will.'' But when you then, on this presupposition (Open Letter,
p.~15), refer expressly to Scripture and declare that ``the more one
strives, seeking nourishment for one's faith, to live oneself into
Scripture, the more it summons us to \emph{thinking} about the
highest objects, a thinking not only of a practical but also of a
theoretical kind''; then I wager ten to one, Doctor, that you are
here confusing, in an altogether unclear way, \emph{religious
thinking}---or, if you will, \emph{religious theory}---with
scientific thinking. Wherever in the world have you come to take the
``learning,'' the ``knowledge'' and the ``wisdom'' of which Christ
and the Apostles speak for scientific knowledge? Has Christ, then,
really the character of a man of science? Have isagogics,
hermeneutics, exegetics, apologetics and dogmatics perhaps, by united
endeavours, at last brought it so far that Christ's words, cited by
you yourself (Matthew 11, 25: ``I praise thee, Father, that thou
hast hid these things from the wise and prudent, and hast revealed
them unto babes''), can now really be taken as a kind of scientific
introduction---an introduction to the science of miracles, the
science that will kindly explain the miracles to us? Or how has it
been possible for you, in one and the same breath, to cite the two
passages of Paul, first 1 Cor. 1, 14,%
% So printed ("1. Cor. 1, 14"); the words quoted are 1 Cor. 1:19.
{} ``the word of the Lord by the prophet: I will cast away the wisdom
of the wise, and bring to nothing the understanding of the
prudent,'' and then the prayer for the Colossians, ``that they might
be filled with the cognition of God's will in all wisdom and
spiritual understanding'' (Col. 1, 9), and here think of scientific
wisdom, scientific understanding? I see well enough that you have
emphasized the words \emph{all wisdom and spiritual understanding};
but I see also that you have forgotten to emphasize the words
\emph{the cognition of God's will}. Had you noticed that the
cognition of which Paul
% --- p. 15 ---
\opage{15}%
speaks is precisely the cognition \emph{of God's will}, you could not
possibly have overlooked that this cognition of the will, for all
theory, is and remains of a practical kind.

You say (Open Letter, p.~17): ``Besides the testimony of the Church
and of Scripture one knows, after all, that of the \emph{Spirit},
the testimony of the Holy Ghost. Should there not be a
\emph{scientific} expression for this testimony, so that one might
have in it a principle well in accord with the Church and Scripture,
yet not drawn from their testimony?''\footnote{From Grundideernes
Logik I (pp.~48--53; 239--43) and II (pp.~151--57; 343--47) it may be
seen that from the standpoint of knowledge I have by no means been
without attention to what the theologians call \textit{testimonium
spiritus}.} This expectation will be disappointed. It is, after all,
a fact that the Greek-Catholic theologian has appealed to the
testimony of the Holy Ghost just as much as the Roman-Catholic, the
Reformed theologian just as much as the Lutheran; and what has been
the consequence? That the one, by virtue of ``the testimony of the
Spirit,'' has pronounced a curse upon the other! So it is when we
look at the history of the Church. If we look at the origin, then
the Holy Ghost gives itself its very mightiest testimony in
Inspiration. What, now, according to the Apostle Peter's speech
(Acts 2, 14--32), is this testimony? ``And it shall come to pass in
the last days, saith God, I will pour out of my Spirit upon all
flesh: and your sons and your daughters shall prophesy, and your
young men shall see visions, and your old men shall dream dreams,''
and so on; but such a testimony does not point in a scientific
direction.

How religious cognition, from the standpoint of religion, will come
to relate itself to scientific cognition is a great and comprehensive
problem, whose solution must be attempted in a philosophy of
religion; but a philosophy of religion that is to solve the task in
any degree must then know how to set the light of knowledge in such
a way that, instead of casting a confusing twilight glimmer over the
miracles, it can illuminate from all sides
% --- p. 16 ---
\opage{16}%
the religious independence of the principle of will, and therewith
of the principle of revelation and of faith.

What I have here, in all brevity, written in reply to your esteemed
Open Letter, I now ask you, Doctor, to receive with goodwill for
closer consideration. We could, as I said, come to agreement on very
essential points, if only we could agree that a science that
explains miracles, a science of miracles, belongs among the
impossible things.

\begin{center}
\rule{0.25\textwidth}{0.4pt}
\end{center}

\end{document}
